%=====================================================================
\chapter{Writing the Thesis and the Paper}\label{ch:writing}
%=====================================================================
\locator{6}{Part~6 --- Producing and Defending Research}

\begin{outcomes}
By the end of this chapter you will be able to:
\begin{tight}
  \item \textbf{State} your contribution in one sentence and build the document
    around it.
  \item \textbf{Structure} an IMRaD document and say what belongs in each part
    and what does not.
  \item \textbf{Write} an abstract in five sentences.
  \item \textbf{Separate} results from interpretation.
  \item \textbf{Assess} your own draft against the reporting criteria you will
    be marked on.
\end{tight}
\end{outcomes}

\begin{prereq}
\chref{synopsis} --- much of the thesis is the synopsis, revised in the past
tense. \chref{validity} for the threats section, and \chref{appraisal} for how
your reader will read you.
\end{prereq}

\begin{keyterms}
IMRaD $\bullet$ storyline $\bullet$ contribution sentence $\bullet$ funnel
introduction $\bullet$ related work $\bullet$ topic sentence $\bullet$ hedging
$\bullet$ signposting $\bullet$ revision
\end{keyterms}

\section{One Sentence First}

Before writing anything, complete this:

\begin{center}
\small
\begin{tabular}{@{}L{14.0cm}@{}}
\toprule
``This work shows that \underline{\hspace{60mm}}, which matters because
\underline{\hspace{45mm}}.'' \\
\bottomrule
\end{tabular}
\end{center}

Everything in the document either supports that sentence or does not belong.
This is the fastest available cure for the two commonest structural problems in
a thesis: a chapter that exists because the work was done rather than because
the argument needs it, and a related-work section that summarises the field
instead of positioning the contribution.

\begin{conceptbox}[Thesis and paper are different documents]
They are not the same content at different lengths.

\begin{center}
\small
\begin{tabular}{@{}L{3.0cm}L{5.4cm}L{5.5cm}@{}}
\toprule
 & \textbf{Thesis} & \textbf{Paper} \\
\midrule
Audience   & Examiners, who assess competence
           & Peers, who want the result \\
Shows      & That you can conduct research independently
           & That the finding is sound and new \\
Method     & Full detail, including what failed
           & Enough to reproduce, compressed \\
Literature & A chapter, comprehensive
           & A section, positioning only \\
Failures   & Reported --- they demonstrate judgement
           & Usually omitted unless they bear on the claim \\
Length     & As long as the argument needs
           & As short as the venue permits \\
\bottomrule
\end{tabular}
\end{center}

\medskip
Write the thesis for the examiner and cut the paper from it, never the reverse:
a thesis assembled from published papers usually has no through-line, and
examiners notice the seams.
\end{conceptbox}

\section{IMRaD, and What Does Not Belong}

\begin{center}
\small
\begin{longtable}{@{}L{2.8cm}L{5.2cm}L{6.2cm}@{}}
\toprule
\textbf{Section} & \textbf{Contains} & \textbf{Does not contain} \\
\midrule
\endfirsthead
\toprule
\textbf{Section} & \textbf{Contains} & \textbf{Does not contain} \\
\midrule
\endhead
\bottomrule
\endfoot
Introduction
  & The problem, the gap, the contribution, the roadmap
  & Method detail; results; a history of the field \\
Related work
  & Positioning against the closest prior work
  & Everything you read; a chronological list \\
Method
  & What you did, in enough detail to repeat
  & Justification of the field's standard practices; results \\
Results
  & What was observed, with uncertainty
  & \textbf{Interpretation}; explanation of why \\
Discussion
  & What it means, limits, comparison to prior work
  & \textbf{Repetition} of the numbers; new results \\
Conclusion
  & The contribution restated, and what follows
  & Anything not already established above \\
\end{longtable}
\end{center}

\begin{alertbox}[The two boundary violations that reviewers always catch]
\textbf{Interpretation in Results.} ``Checklist review found 1.3 more defects
($p = 0.004$), \emph{showing that structured guidance improves reviewer
attention}.'' The italicised clause is a claim about mechanism and belongs in
Discussion. Keeping Results free of interpretation lets a reader who disagrees
with your interpretation still use your data --- which is what makes a result
durable.

\medskip
\textbf{Numbers in Discussion.} A Discussion that restates every result before
interpreting it doubles its length and interprets nothing. Refer to results;
do not repeat them.
\end{alertbox}

\section{The Abstract in Five Sentences}

\begin{center}
\small
\begin{tabular}{@{}L{0.6cm}L{3.4cm}L{9.8cm}@{}}
\toprule
1 & Context and problem & What area, and what is wrong \\
2 & Gap               & What specifically is not known \\
3 & What you did      & Design, population, scale --- one sentence \\
4 & What you found    & The actual result, with a number \\
5 & What it means     & The implication for practice or research \\
\bottomrule
\end{tabular}
\end{center}

\begin{pitfall}[The abstract that promises without delivering]
``\ldots{} The results are presented and discussed, and recommendations are
made.'' Sentence 4 has been replaced by a description of the fact that there is
a sentence 4.

\medskip
A reader deciding whether to read your paper needs the finding, not the
assurance that one exists. Put the number in the abstract. If the result is
null, say so with the interval (\chref{inference}) --- a well-powered null
stated plainly is more informative than a vague positive.
\end{pitfall}

\section{The Introduction Funnel}

\begin{center}
\small
\begin{tabular}{@{}L{3.6cm}L{10.7cm}@{}}
\toprule
Paragraph 1 & The broad area and why it matters. Two or three sentences, not a
              page \\
Paragraph 2 & Narrowing: within this area, this specific problem \\
Paragraph 3 & What is known, briefly, with citations \\
Paragraph 4 & The gap. Specific, evidenced \\
Paragraph 5 & What this work does, and the contribution stated as a list \\
Paragraph 6 & Roadmap, one sentence per section \\
\bottomrule
\end{tabular}
\end{center}

\begin{conceptbox}[State contributions as a list]
Reviewers look for it, and writing it forces precision:

\medskip
``This paper makes three contributions. First, a controlled experiment with 34
professional developers measuring the effect of checklist-guided review under
time-boxing --- the first to vary time budget as a factor. Second, evidence
that the checklist benefit reverses below a 15-minute budget. Third, a
replication package containing the tasks, seeded defects and analysis scripts.''

\medskip
If you cannot write this list, the work is not finished --- or the contribution
is not yet identified (\chref{contribution}).
\end{conceptbox}

\section{Tables and Figures}

\begin{center}
\small
\begin{tabular}{@{}L{3.4cm}L{10.9cm}@{}}
\toprule
Self-contained     & Caption plus content readable without the body text.
                     Readers skim figures first \\
Caption says the finding & ``Defects found by condition and time budget; the
                     checklist advantage reverses below 15 minutes'' beats
                     ``Results by condition'' \\
Show uncertainty   & Intervals or distributions, not bare bars
                     (\chref{inference}) \\
Show the data      & Where $n$ is small, plot the points. A bar chart of two
                     means hides everything \\
Table or figure, not both & Never present the same data twice \\
Accessible         & Distinguishable without colour; readable when printed
                     greyscale, which is how a thesis is often examined \\
\bottomrule
\end{tabular}
\end{center}

\section{Paragraphs}

\begin{center}
\small
\begin{tabular}{@{}L{3.4cm}L{10.9cm}@{}}
\toprule
One idea per paragraph & If it has two, it is two paragraphs \\
Topic sentence first   & The first sentence states the paragraph's claim; the
                         rest supports it. A reader should be able to read only
                         the first sentences and follow the argument \\
Old before new         & Begin with what the reader already knows, end with
                         what is new. The new thing becomes the old thing for
                         the next sentence \\
Hedge accurately       & ``May suggest a possible tendency'' hedges three
                         times. Hedge once, precisely: ``is consistent with'',
                         ``suggests'' \\
Cut throat-clearing    & ``It is important to note that'' adds nothing. Delete
                         and start with the noun \\
\bottomrule
\end{tabular}
\end{center}

\begin{conceptbox}[The first-sentence test]
Read only the first sentence of every paragraph in a section, in order. Does
the argument hold together? Can you tell what the section claims?

\medskip
If not, the paragraphs are organised by what you did rather than by what you
are arguing. This test takes five minutes per chapter and finds more structural
problems than any amount of re-reading.
\end{conceptbox}

\begin{alertbox}[Writing in a second language]
Most scholars here write in English as a second or third language, and the
disadvantage is real. Three things help more than general advice to ``improve
your English''.

\medskip
\textbf{Read for structure, not vocabulary.} Take a well-written paper in your
area and mark its topic sentences. The rhetorical moves --- how a gap is set
up, how a limitation is conceded --- are learnable patterns, and they matter far
more than vocabulary range.

\medskip
\textbf{Write simple sentences.} Short declarative sentences are clearer for
every reader and much harder to get wrong. Complex subordination is where
second-language writing usually breaks, and it buys nothing.

\medskip
\textbf{Get a human reader.} Your supervisor, a writing centre, a colleague.
Note that \chref{genai} constrains what tooling you may use here, and that
constraint is stricter than international norms --- so build the human route
early rather than relying on something you may not use.
\end{alertbox}

\section{The Standard You Are Marked Against}

\begin{regbox}[Quality in reporting --- PhD \S14.2]
``The quality of presentation and reporting in the dissertation shall reflect
the following characteristics:
\begin{tight}
  \item[(i)] The document should be well written.
  \item[(ii)] The contents should be balanced, well organized, appropriately
    styled; clearly structured, and well cohered.
  \item[(iii)] The document should [be] free from grammatical and spelling
    errors and flawed terminology.
  \item[(iv)] Minor shortcomings such as inaccurate use of acronyms and
    clumsy-looking sentence structure have been addressed.
  \item[(v)] Quantitative research proposals must include a valid statistical
    design for data analysis.
  \item[(vi)] Formatting shall be according to the IUB thesis guidelines
    standards.''
\end{tight}
\end{regbox}

\begin{conceptbox}[Use it as a checklist, because it is one]
Item (vi) is the one that causes avoidable returns: the
\emph{IUB Theses/Synopses Guidelines-2023}~\cite{iub2023theses} specify
formatting, and a thesis that deviates is sent back before anyone reads the
research. Get the template at the start and write in it.

\medskip
Item (v) is why the analysis plan belongs in the synopsis and not in the
write-up: a valid statistical design is required in the \emph{proposal}, which
means it must exist before the data does (\chref{prereg}).
\end{conceptbox}

\section{Revision}

\begin{center}
\small
\begin{tabular}{@{}L{3.0cm}L{11.3cm}@{}}
\toprule
Pass 1: structure & Does each chapter earn its place? Does the first-sentence
                    test pass? Ignore wording entirely \\
Pass 2: argument  & Is every claim supported? Is every number traceable to the
                    analysis? Is anything claimed that the design cannot
                    support (\chref{validity})? \\
Pass 3: paragraph & Topic sentences, old-before-new, one idea each \\
Pass 4: sentence  & Length, hedging, throat-clearing, terminology consistency \\
Pass 5: mechanics & References, cross-references, figure numbers, formatting
                    against the template, spelling \\
\bottomrule
\end{tabular}
\end{center}

Do these as separate passes. Fixing commas while deciding whether a chapter
belongs is how drafts stall: the small work is comfortable and the large work
is not.

\begin{traditions}{writing it up}
\tradEMP{IMRaD with a rigid Results/Discussion boundary, effect sizes with
intervals, and a threats section that grades each threat (\chref{validity}).}
\tradDSR{Problem, objectives, artefact, evaluation, design knowledge. The
design knowledge must be stated as a separate claim, not left implicit in the
artefact description (\chref{designscience}).}
\tradQUAL{Thick description, quotations as evidence with attribution, and the
analytic ladder from data to interpretation (\chref{qda}). Longer, and the
length is required by the method.}
\tradFORM{Definitions, theorem, proof, discussion of assumptions. The
assumptions section is where a formal paper is most often weak and most often
attacked.}
\end{traditions}

\begin{lab}[Test a draft chapter]
\begin{tightnum}
  \item Write your contribution sentence. Check every section of the chapter
    against it; mark anything that does not support it.
  \item Extract the first sentence of every paragraph and read them in order.
    Fix what does not follow.
  \item Find every interpretive claim in your Results section and move it.
  \item Rewrite your abstract in exactly five sentences, with a number in
    sentence 4.
  \item Check the chapter against \S14.2, item by item.
\end{tightnum}
\end{lab}

\begin{checkpoint}
\begin{tightnum}
  \item Why should a thesis be written first and papers cut from it?
  \item Give an example of a sentence that belongs in Discussion but is usually
    written in Results.
  \item What does the first-sentence test detect that ordinary re-reading does
    not?
  \item Which clause of \S14.2 most often causes a thesis to be returned before
    it is read?
\end{tightnum}
\end{checkpoint}

\begin{chaptersummary}
\begin{tight}
  \item Write the contribution sentence first. Everything either supports it or
    is cut.
  \item Thesis and paper have different audiences and different content. Write
    the thesis, cut the paper.
  \item Keep Results free of interpretation and Discussion free of repetition.
  \item Five-sentence abstract, with an actual number in sentence 4.
  \item State contributions as an explicit list.
  \item Captions state findings; figures show uncertainty and, at small $n$,
    the data.
  \item Topic sentence first, old before new, hedge once. Apply the
    first-sentence test to every section.
  \item \S14.2 is the marking criterion and \S14.2(vi) means writing in the IUB
    template from the start.
  \item Revise in five separate passes, largest concern first.
\end{tight}
\end{chaptersummary}

\begin{reviewq}
  \item The Results/Discussion boundary is a convention. What epistemic purpose
    does it serve, and are there computing subfields where it should be
    abandoned?
  \item Second-language writers face a real disadvantage in a system where
    English is the language of publication, and \chref{genai} restricts the
    tooling that might offset it. What should an institution provide instead?
  \item A thesis by publication --- stapled papers with a linking chapter --- is
    permitted at some universities. Argue for and against, with reference to
    what \S14.2(ii) means by ``well cohered''.
  \item \S14.2(v) requires a valid statistical design in the \emph{proposal}.
    What would it take to enforce this meaningfully at the DRC stage, and what
    would change if it were enforced?
\end{reviewq}
